\section{Architecture: a scoped obligation graph around a shared fact store}
\label{sec:architecture}
\status{Specified architecture. The Python prototypes implement several workers
and certificate formats; none of them implements this Lean controller.}

\subsection{Three data structures that must never be confused}

A single global collection of ``facts'' is too ambiguous to be safe. The
controller maintains three stores with different admission rules.

The \emph{proved-fact store} contains propositions together with Lean proof
expressions valid in a particular local context. Each entry records its
source, its assumptions, and the proof dependencies needed to reconstruct it.
A worker may consume it as evidence.

The \emph{obligation graph} contains goals waiting for proofs: ordinary goals,
side conditions, candidate lemmas, induction cases, witness specifications,
and proof-reconstruction holes. An AND node is solved only when all of its
required children are solved; an OR node records alternative approaches to the
same goal. Failed attempts and exhausted budgets do not become logical facts.

The \emph{search store} contains models, candidate terms, speculative
rewrites, usefulness scores, and unsuccessful attempts. These may direct
search. They are never accepted as hypotheses by a proof worker. In
particular, a conjectural equality may merge nodes in an isolated exploration
graph, but it must not merge nodes in the proved equality state.

\begin{invariant}[The ledger invariant]
Every entry in the proved-fact store carries a proof term and a scope. Nothing
enters it that is not accompanied by both. Membership of the proved-fact store
is the \emph{only} licence a worker has to use a proposition as a premise.
\end{invariant}

At the outer layer these stores are organised as an AND/OR search graph over
structural choices; at each open leaf a coordination layer exposes certified
facts and typed views to specialised workers. \grind{} maintains its own
efficient internal equality and theory state behind that boundary. The
acceptance layer reconstructs a Lean proof of each returned proposition,
checks scope and dependencies, and ultimately constructs a proof of the
original goal. Numerical success, satisfiability status, an unverified
certificate, and a plausible conjecture are never accepted as facts.

\subsection{The demand protocol and the worker contract}

Workers are addressed through an explicit vocabulary of demands rather than
being invoked as opaque tactics. The following is a proposed interface, not a
transcription of a released Lean API.

\begin{lstlisting}
Demand :=
  | Prove          (goal)
  | ProveGuard     (context, proposition)
  | FindWitness    (context, binders, specification, grammar)
  | Instantiate    (context, theorem, desired_consequence)
  | Generalize     (context, recursive_call, blocked_goal)
  | FindLemma      (context, left_frontier, right_frontier)
  | ProveInduction (context, motive, recursor_candidates)
  | Reconstruct    (context, source_goal, external_certificate)

Worker.supports : GoalView -> SupportReport
Worker.step     : Snapshot -> DeltaFacts -> Budget -> WorkerResult

WorkerResult :=
  | Proved     (lean_proof : CheckedProofRef)
               (dependencies) (replay_record)
  | Progress   (facts       : Array ScopedFact)
               (obligations : Array ConditionalObligation)
               (continuation : OpaqueToken)
  | Candidate  (kind : CertificateKind) (payload : BoundedData)
               (proposed_next_demands)
  | Unknown    (reason : FailureReason) (diagnostics) (continuation)

ScopedFact := {
  proposition  : Expr,
  proof        : Expr,
  scope        : ScopeId,
  dependencies : Array FactId
}
\end{lstlisting}

\code{CheckedProofRef} is a capability issued by the acceptance layer, never by
an external process. An external solver can return only data or a candidate. A
native Lean tactic worker may return a proof expression, but the coordinator
still verifies its type, its local-variable support, its auxiliary
declarations, and its allowed axioms before committing it.

A worker returning \code{Progress} must not pretend that its residual goals are
solved. The controller may tentatively construct a parent proof containing
metavariables, but overall success is reported only after those metavariables
are assigned and the completed expression is checked.

A \emph{conditional obligation} is a directed dependency, not an assertion.
Cancelling $a$ in a field creates the obligation $a\ne0$; until that proof
exists, the cancelled equality cannot enter the fact store. The result protocol
must keep a conjectured fact visibly different from a proved fact even when the
two share the same proposition.

A counterexample belongs in \code{Candidate} unless its terms, its hypotheses,
and the negation of the source conclusion have all been validated in Lean. For
an abstraction-based translation, a satisfying assignment may be no more than a
model of an overapproximation.

\subsection{Three outcomes, not an overloaded Boolean}

The tactic's result type should distinguish \emph{proved}, \emph{refuted with
a checked countermodel}, and \emph{unknown}, and the diagnostic attached to
\unknown{} should name the first substantive blocked obligation or the
exhausted budget. Collapsing these into success/failure destroys exactly the
information the next attempt needs, and it invites the most common category
error in this area: reading a failed bounded search as a disproof.

\begin{remark}
Throughout this document, a search failure is \unknown{}. It is never evidence
that the requested theorem is false.
\end{remark}

\subsection{Snapshots, scope keys, and proof transport}

A proposed \code{GoalSnapshot} contains the environment revision, the local
context, local instances, the metavariable context, the target expression, the
transparency configuration, branch assumptions, the normalisation version, and
worker state. The controller must restore \emph{all} of it on backtracking.
Restoring only the visible list of hypotheses is insufficient: a worker may
have instantiated metavariables or created local declarations whose identities
are used by other expressions.

A persistent fact key contains the expression, its type, universe levels where
relevant, the identity of the ordered ring or field instance, a local-context
fingerprint, and a scope identifier. Printed text is not a safe key --- names
printed as \code{x} can denote different free variables in sibling goals ---
and expression pointers are not suitable for caches that outlive an
environment. Implicit typeclass arguments matter: two syntactically similar
applications of an overloaded operation can mean different things because their
instances differ.

Typed views must therefore preserve source semantics, and the traps are
concrete:

\begin{itemize}[leftmargin=1.4em]
\item $x-y$ over $\N$ is truncated subtraction; over $\Z$ or $\R$ it is ring
  subtraction. A translator that treats these identically is unsound.
\item A \code{Fin n} expression is not automatically an integer with
  unrestricted arithmetic.
\item A proposition about real division cannot be transported through a
  denominator without a nonzero or sign guard. Clearing denominators in an
  inequality requires a sign proof. These become \code{ProveGuard}
  obligations; if the guard fails, the rewrite is unavailable, not
  approximately valid.
\item An integer embedding into a field can be injective; a fixed-width
  bit-vector projection is not. Mixing those relationships inside a generic
  ``numeric cast'' cache would be a serious error.
\end{itemize}

Facts proved using a branch hypothesis can be shared only within that branch or
a descendant. To move one upward, discharge the hypothesis and transport an
implication, or combine proofs from a complete case split. This applies as
strongly to numerical bounds, inferred equalities, and cached certificates as
it does to ordinary theorem applications. There is no ``most likely branch''
rule in the logical layer.

A subtler issue is proof transport. A successful worker may introduce auxiliary
theorem declarations and return a proof referring to them; restoring the
environment while retaining only the final expression would leave dangling
references. A commit must preserve the required declaration closure, after
checking it, or reconstruct a self-contained proof term in the parent
environment. One safe reuse method abstracts a proof over its exact
free-variable dependencies and then instantiates that closed theorem in the
destination context. External results must be re-elaborated against the exact
destination context rather than transporting free-variable identifiers between
unrelated goals.

\subsection{Facts carry explanations, not only truth flags}

A nonlinear worker may prove $0<c$ from polynomial constraints. That proof
justifies $c\ne0$, which unlocks a division rewrite; the rewrite's proof yields
an equality; the equality may enable congruence closure; and that may make an
E-matching instance useful. The chain should be represented as a proof DAG with
explicit edges, not rediscovered by reparsing tactic text.

This is why an arithmetic worker returns individual proved consequences and
their assumptions rather than a Boolean ``solved'' flag. A checked sign or
equality is useful well beyond the goal that prompted it.

\subsection{Event-driven cooperation}

Workers subscribe to deltas rather than repeatedly scanning the entire context.
Typical events are a new bound on an arithmetic expression, a new equality
joining two term classes, a proved side condition, an instantiated theorem, and
a changed target after structural decomposition.

The coordinator normalises a returned fact, checks whether it adds information,
stores its proof, and notifies only relevant subscribers. Changing an interval
bound invalidates or strengthens a Bernstein search problem; learning a
polynomial equality may reduce the dimension of a subsequent positivity
certificate; proving an existential witness's denominator nonzero activates its
algebraic replay.

Logical circularity and scheduling cycles must be distinguished. A worker
waiting for a fact that another worker might establish is an ordinary
dependency cycle; it may call for a different search action. It does not permit
both workers to assume their own outputs. Track these dependencies as strongly
connected components and either pursue an independent obligation, perform a
justified case split, or return a diagnostic explaining the blockage.

\subsection{An explicit scheduler}

The initial implementation should use deterministic, cost-aware scheduling and
must not require a trained model. Each task has a relevance estimate, an
expected information gain, a recent cost, an age, and a resource ceiling. A
simple score is
\[
  s(a)=w_r R(a)+w_g G(a)+w_a A(a)-w_c\widehat C(a),
\]
with the weights fixed in the experiment manifest. This is a heuristic, and
the important part is not its shape but two disciplines around it.

First, \emph{reserve a minimum quota for every eligible worker class}, so that
a high-scoring stream of cheap actions cannot starve structural search. A cheap
task that regenerates itself is the normal failure mode here.

Second, \emph{the estimated benefit of an action must include the cost of the
certificate it will produce}, not only the probability that it closes the goal.
A worker that reliably succeeds while emitting a proof too large to elaborate
is not a good scheduling choice, and a scheduler that scores only success
probability will keep choosing it.

\begin{lstlisting}
solve(goal, mode, budgets):
  if mode == compatibility:
    run exact_stock_grind(goal, budgets.stock)
    if accepted: return proof
    restore_original_goal()
  agenda := initial_structural_and_local_tasks(goal)
  while resources_remain:
    action   := choose_with_reserved_worker_quotas(agenda)
    snapshot := save_full_branch_state()
    result   := run_bounded(action)
    restore_or_commit_only_checked_changes(snapshot, result)
    if original_goal_has_checked_proof(): return proof
    enqueue_new_obligations_and_relevant_delta_subscribers()
  return Unknown(with_structured_diagnostics)
\end{lstlisting}

Heartbeats alone are insufficient for external numerical code. Maintain
separate limits on wall time, memory, instantiated theorems, polynomial
monomials, rational bit lengths, certificate size, induction depth, and replay
work. External workers should run in bounded child processes without network
access. A process timeout returns \unknown{} and cannot create a theorem.

Finite budgets imply incompleteness. Fair worker quotas do not make
unrestricted Lean proof search complete; they make the search policy
interpretable and reduce avoidable starvation. A deterministic mode, in which
the schedule is a function of the goal and the manifest alone, should be
available for reproducible experiments even if an interactive mode is allowed
to be adaptive.
